\chapter{Lanskap Regulasi dan Kewajiban Kepatuhan}

\section*{Tujuan Pembelajaran}
\begin{itemize}
  \item Mengidentifikasi regulasi dan standar yang mengikat sebuah organisasi berdasarkan sektor, jenis data, dan model bisnisnya.
  \item Menurunkan pasal regulasi menjadi kewajiban yang dapat diperiksa, lengkap dengan pemilik dan buktinya.
  \item Membedakan kewajiban hukum, kewajiban kontrak, dan standar sukarela beserta akibat pelanggarannya.
\end{itemize}

\section{Pendahuluan}

% TODO
% Cakupan: aturan baru terbit, tetapi tidak ada yang tahu siapa yang harus bertindak; kepatuhan baru dikerjakan menjelang audit.

\section{Membaca Regulasi Indonesia}

% TODO
% Cakupan: hierarki peraturan perundang-undangan menurut UU 12/2011 beserta perubahannya; beda UU, PP, Perpres, peraturan menteri, peraturan lembaga, dan surat edaran; cara memeriksa status berlaku di JDIH BPK dan peraturan.go.id.

\section{Pelindungan Data Pribadi dan Sistem Elektronik}

% TODO
% Cakupan: UU 27/2022 dan PP 33/2026 dalam garis besar (rincian di Bab 7); UU ITE beserta perubahan terakhir melalui UU 1/2024; PP 71/2019 tentang PSTE: PSE lingkup publik dan privat, pendaftaran, keandalan dan keamanan sistem.

\section{Regulasi Sektor Kesehatan}

% TODO
% Cakupan: UU 17/2023 tentang Kesehatan dan PP 28/2024; Permenkes 24/2022 tentang Rekam Medis: kewajiban rekam medis elektronik, masa simpan, kerahasiaan, dan interoperabilitas dengan SATUSEHAT; aturan Kemenkes lain yang relevan (periksa status terbaru saat menulis isi).
% Kasus Medikara: integrasi point-to-point ke SATUSEHAT dan BPJS Kesehatan serta SIMRS yang berbeda versi di tiap rumah sakit.

\section{Keamanan Siber Nasional dan Infrastruktur Informasi Vital}

% TODO
% Cakupan: peran BSSN; Perpres 82/2022 tentang Pelindungan Infrastruktur Informasi Vital dan posisi sektor kesehatan; Peraturan BSSN 1/2024 tentang Pengelolaan Insiden Siber; status RUU Keamanan dan Ketahanan Siber (periksa ulang).

\section{Regulasi yang Mengikat Mitra: OJK dan SPBE}

% TODO
% Cakupan: POJK 4/POJK.05/2021 bagi mitra asuransi sebagai lembaga jasa keuangan nonbank; kewajiban mitra yang mengalir ke Medikara lewat kontrak; Perpres 95/2018 dan evaluasi SPBE bagi mitra rumah sakit pemerintah (arsitektur SPBE dibahas rinci di mata kuliah IT40513).
% Kasus Medikara: mitra asuransi dan layanan kesehatan karyawan perusahaan dalam rencana ekosistem.

\section{Register Kewajiban Kepatuhan}

% TODO
% Cakupan: struktur register (sumber, pasal, kewajiban, pemilik, bukti, status); cara menurunkan pasal menjadi kewajiban yang dapat diperiksa; prinsip ISO 37301 tentang sistem manajemen kepatuhan; beda kewajiban hukum, kontrak, dan sukarela.
% Pendalaman teknis: menautkan kewajiban ke persyaratan sistem yang dapat diuji (logging, enkripsi, masa simpan) dan ke backlog tim pengembang.

\section{Aktivitas: Register Kewajiban Kepatuhan}

% TODO
% Tujuan aktivitas: memiliki satu daftar kewajiban yang jelas pemilik dan buktinya, sebagai dasar seluruh pekerjaan kepatuhan semester ini.
% Peran dalam tim: analis regulasi memimpin pembacaan pasal (anggota berlatar hukum sangat membantu); pemilik konteks menunjuk unit pemilik; analis teknologi menandai kewajiban yang berdampak pada sistem; analis risiko menilai akibat bila kewajiban dilanggar.
% Langkah: tentukan sektor, jenis data, dan model bisnis; daftar regulasi yang berlaku; turunkan pasal penting menjadi kewajiban; isi pemilik, bukti, dan status awal; tandai kewajiban berdampak sistem.
% Luaran dan tonggak capstone: register kewajiban kepatuhan; masuk C1 dan dipakai ulang pada Bab 5, 7, dan 12.
% Berkas sources yang direncanakan: sources/chapter03/register-kewajiban.xlsx, sources/chapter03/daftar-regulasi-medikara.md

\section{Ringkasan}

% TODO

\section{Latihan}

\begin{enumerate}
  \item Sebuah klinik mitra menyimpan rekam medis di layanan cloud luar negeri. Kewajiban apa saja yang perlu diperiksa, dan pada peraturan mana?
  \item Tulis ulang satu pasal UU PDP menjadi kewajiban yang dapat diperiksa auditor: siapa melakukan apa, kapan, dan apa buktinya.
  \item Jelaskan mengapa sertifikat ISO/IEC 27001 tidak otomatis membuktikan kepatuhan pada UU PDP.
\end{enumerate}
